% Soil zone tests (T1-T7, inputs, robustness) and regression tests
\section{Soil zone: the column, the hillslope and the interface}\label{sec:vsoilzone1}
The first soil-zone tests (\cmd{tests/soilzone}) put the water table into the soil zone of small synthetic models
(\cmd{sz\_models.py}: a generated MODFLOW model under Raven HRUs with soil, forcing written by the script) and compare
them with exact values or with a MODFLOW-only reference in which the soil zone is a MODFLOW layer and a drain at the land
surface takes the overflow (\cmd{sz\_reference.py}). Every coupled run is also checked by \cmd{sz\_check.py}, which
closes, step by step, MODFLOW's balance, Raven's balance and the recharge, soil-zone and unsaturated-zone identities of
Section~\ref{sec:accsz2} at their standard limits (recharge $5\times10^{-9}$, soil zone and columns $10^{-8}$ relative;
\cmd{sz\_check.ID\_LIMITS}), and requires no soil-zone shortfall. Unlike the coupled tests of the earlier sections these
tests do not start from the Liard examples (except T7). The numbers are those of the full mode of
\cmd{tests/run\_all.py} with MODFLOW~6.8.1.

\begin{longtable}{P{0.13\linewidth}P{0.42\linewidth}P{0.36\linewidth}}
\caption{Soil zone: the column, the hillslope, time steps, restarts and the Liard model (\cmd{test\_column},
\cmd{test\_hillslope}, \cmd{test\_dynamics}, \cmd{test\_restart}, \cmd{test\_liard\_physical}, \cmd{soilzone\_unit},
\cmd{spike\_xmi}).}\label{tab:vsoilzone1}\\
\toprule\hdr Test & What is checked & Result\\\midrule\endfirsthead
\caption[]{Soil zone (continued).}\\
\toprule\hdr Test & What is checked & Result\\\midrule\endhead
Unit & \cmd{soilzone\_unit.cpp}: filling and draining of the layers, the sub-grid level and its inverse, the contact weights and flux derivatives against finite differences, continuity across the sample bases, the local solve (monotone, root, mass), the Robin slope; 20\,000 random cases each & 57 checks, 0 failed\\
Spike & \cmd{spike\_xmi.py}: values written between MODFLOW's outer iterations through XMI (GHB \cmd{BHEAD}/\cmd{COND}, RCH \cmd{RECHARGE}) reach the solver; the coefficients are recomputed every outer iteration; a solve after convergence is a normal iteration & all three hold (MODFLOW 6.8.1)\\
T1 flat column & closed column, $S_y$ = drainable porosity = 0.2, recharge 10~mm/d: the water table rises 50~mm/d below the soil base and inside the soil; water table and soil level agree once connected; the soil zone (400~mm of capacity) is full on day~41 & rise 49.95~mm/d below and inside; largest difference $1.3\times10^{-4}$~m; full on day 41\\
T1b & the same with $S_y=0.08$: 125~mm/d below the base, 50~mm/d inside, storage-consistency warning & 124.7 and 49.95~mm/d; warning given\\
T2 hillslope, wetting & ten 100~m HRUs on a 0.5\,\% slope over a closed 20~m aquifer, rain 10~mm/d for 150 days (Figure~\ref{fig:szhillslope}): against the MODFLOW-only reference & heads within 3.3~cm (mean 0.67~cm), soil level within 3.9~cm, total outflow within $2.5\times10^{-5}$; overflow begins on day 16 in both; the toe connects first (day 7), HRU 2 last (day 73), HRU 1 never\\
T3 hillslope, drying & the 150 dry days that follow: HRUs that connected and whose reference head ends below their base lose their saturation; those whose head stays above it stay saturated; no overflow where the water table is below the land; the upslope overflow stops with the rain; the toe's seepage recedes; \cmd{--selftest} shows that each check fails on altered results & HRUs 2--7: $F=0$ at the end; 8--10: $F=1$; upslope overflow ends on day 150; toe 814~m$^3$/d on day 161, 173.5 on day 300 (reference 173.6); saturated HRUs 9 $\to$ 3\\
T4 time step & the hillslope with steps of 1~d, 6~h and 1~h, each against the reference with the same step: the splitting error of Raven's processes ahead of the exchange shrinks with the step & mean head difference 0.67, 0.24 and 0.14~cm; outflow within $6\times10^{-5}$\\
T5 stiff contact & a thin, conductive soil--aquifer contact ($\dt/t_c\approx48$ for daily steps): the iterated interface against one explicit pass & iterated: at most 9 updates, heads within 4.2~cm of the reference, outflow $8\times10^{-5}$; one pass: 0.99~m and 1.3\,\%\\
SG sub-grid & \cmd{:SoilZoneSubgrid 5} on the same slope: the saturated fraction changes smoothly instead of cell by cell & largest daily change of $F$ 0.42 (1.0 with one point per cell); outflow within 0.4\,\% of the one-point run\\
Restart & 300 days against 150 + 150 days from \cmd{solution.rvc} (\cmd{:GWSoilZoneColumns}); R2: an existing model with a specified-head toe and unsaturated columns, 10 days against 5 + 5 & 1\,500 rows: water table and soil level $2\times10^{-10}$, $F$ $1.2\times10^{-7}$, exchange $2.4\times10^{-6}$ (limit $10^{-4}$); R2 $2.5\times10^{-7}$; R3 (no \cmd{:SoilZone PHYSICAL}, seepage holding the lowest soil layer above capacity, 30 days against 12 + 18): soil storage identical, recharge and seepage within $3.4\times10^{-9}$; R4 (\cmd{KINEMATIC} and \cmd{LAG} with $S_y$ equal to the porosity, 20 days against 10 + 10): column storage as written, restart identical to the continuous run\\
T6 & planned: the Sagehen model of GSFLOW as an existing model with a soil zone (qualitative comparison) & not implemented; the soil zone with existing models is verified by X1--X8 (Table~\ref{tab:vsoilzone}) and the rough-terrain suite (Section~\ref{sec:vrough})\\
T7 Liard & the Liard example with \cmd{:SoilZone PHYSICAL SOIL[0] SOIL[1]} (85 coupled HRUs, 2~km grid, one year from 1985-10-01; soils of drainable porosity 0.02 against $S_y=0.25$, a stiff interface, storage warning expected) & 365 steps, 0 unconverged (MODFLOW and interface), 3\,509 outer iterations (at most 72 in a step), residual $8\times10^{-12}$ of the gross flow, no shortfall; 80 of 85 HRUs more than 20\,\% saturated at some time\\\bottomrule
\end{longtable}

\begin{longtable}{P{0.07\linewidth}P{0.55\linewidth}P{0.30\linewidth}}
\caption{Soil zone: input handling (\cmd{test\_inputs.py}; short runs of the T1 column unless stated).}\label{tab:vszinputs}\\
\toprule\hdr Test & Case & Result\\\midrule\endfirsthead
\caption[]{Soil zone: input handling (continued).}\\
\toprule\hdr Test & Case & Result\\\midrule\endhead
I1 & \cmd{HYDRAUL\_COND} missing for the soil class of the bottom soil-zone layer & stops (``needs HYDRAUL\_COND'')\\
I2 & \cmd{:SoilZoneIteration} with the update count only (0) & accepted; explicit, no interface update\\
I3 & \cmd{:SoilZone PHYSICAL} followed by \cmd{:SoilZone CONCEPTUAL} & the last counts: a conceptual run\\
I4 & soil-zone options under \cmd{CONCEPTUAL} & reported as ignored\\
I5 & a \cmd{:GWSoilZoneColumns} block with a column missing & reported (``do not match its soil-zone columns'')\\
I6 & the same block with its rows in another order & used (matched by cell)\\
I7 & groundwater blocks in the \cmd{.rvc} of a run without \cmd{:GroundwaterModel} & skipped; Raven's initial conditions after them apply\\
I8 & \cmd{:SurfaceWaterExchange} on a soil-zone layer & refused\\
I9 & a damaged \cmd{:GWHeads} header (\cmd{-1 2 2}; \cmd{100000 100000 100000}) & stops with a message\\
I10 & \cmd{:StartTime} in the \cmd{.rvc}, run under valgrind & no uninitialised read in \cmd{ParseInitialConditionsFile} or \cmd{TimeDifference}\\
I11 & \cmd{:BoundaryHead NOPE} beside boundary \cmd{EDGE}; \cmd{:WellFlowRate 7} with only well 1 & stop naming the series and the boundaries or wells there are; the right names run\\
I12 & \cmd{:SeepageToRaven}, \cmd{:KeepPackages}, \cmd{:MF6StartDate}, \cmd{:MF6Model} on a model Raven builds; depth classes and initial flux with \cmd{NONE}; \cmd{:SoilZoneAquiferClass} on a generated model; an unknown profile attribute; \cmd{LAYERS BOTTOM} & the four stop; the next three are reported as ignored; the last two stop\\
I13 & a missing \cmd{.rvg} file; a short \cmd{:Wells} row & named in \cmd{Raven\_errors.txt}; the message names \cmd{:Wells}\\
I14 & age tracer (no groundwater): 20 days against 7 + 7 + 6 days chained through \cmd{solution.rvc} & ages of every store and the outflow within $10^{-4}$~d of the continuous run; \cmd{:StartTime} the start of each part, \cmd{:AgeTracerDatum} 2000-01-01 in the restarted parts\\
I15 & age-tracer chains across year boundaries: 1999-12-25 (7 + 7 + 6 days), 1998-12-20 (12 + 400 + 8 days; full mode 400 + 380) & every part runs; ages within $10^{-4}$~d of the continuous run; \cmd{:AgeTracerDatum} 1999-12-25 and 1998-12-20\\
I16 & age tracer started from the \cmd{solution.rvc} of a 30-day spin-up without it & ages 0 at the start, 1~d after the first day; no \cmd{:AgeTracerDatum}\\
I17 & EnKF, 3 members, one well, \cmd{:ExtraRVTFilename} with $-100$, $-200$, $-300$~m$^3$/d & each member pumps its own rate; a series naming no well stops the run\\
I18 & \cmd{:GWSoilZoneColumns} with \cmd{nan} or $-50$~mm, \cmd{:GWHeads} with \cmd{nan} & each stops with a message naming the block; the undamaged file runs\\
I19 & \cmd{:Transport Nitrate} (\cmd{CONCEPTUAL}, \cmd{PHYSICAL}) and \cmd{:Transport AGE} with \cmd{:GroundwaterModel MODFLOW6}; the same model without groundwater & the coupled runs stop naming the constituent; the run without groundwater completes\\\bottomrule
\end{longtable}

\begin{longtable}{P{0.07\linewidth}P{0.55\linewidth}P{0.30\linewidth}}
\caption{Soil zone: robustness (\cmd{test\_robustness.py}; cases found by the code inspections). R14--R20 are the negative-recharge
cases N1--N7 of Section~\ref{sec:vsoilzone3integ}.}\label{tab:vszrobust}\\
\toprule\hdr Test & Case & Result\\\midrule\endfirsthead
\caption[]{Soil zone: robustness (continued).}\\
\toprule\hdr Test & Case & Result\\\midrule\endhead
R1 & stiff flat column: drainable porosity 0.02 against $S_y=0.25$, high conductance & every step converged; at most 12 updates, 27 outer iterations\\
R2, R2c & one HRU over ten cells against ten one-cell HRUs; R2c with constant-rate percolation & outflow $-0.2\,\%$, MODFLOW recharge ratio 1.00; R2c outflow $-0.2\,\%$\\
R3 & \cmd{RECHARGE\_DELAY} with \cmd{:SoilZone PHYSICAL} & refused\\
R4 & a smoothing depth below 0.1~m & warned\\
R5 & hourly steps, smoothing depth 0.001~m & no unconverged step; MODFLOW error $1.1\times10^{-11}$\\
R6 & slope 0.1 with $5\times5$ DEM points per cell & converged; at most 8 updates\\
R7 & explicit mode (\cmd{:SoilZoneIteration 0}) & no shortfall\\
R8 & one column $10^{-6}$~mm wetter than the others & heads stay uniform ($3.7\times10^{-9}$~m)\\
R9 & hotstart columns that cannot hold the \cmd{.rvc} layer split & 4\,000~m$^3$ redistribution reported in step 1, 0 later\\
R10 & water rising into a thin soil on a steep, low-storage slope, update budget 1 and explicit mode & no shortfall\\
R11 & explicit mode on the R2 hillslope against the iterated run & outflow within 0.6\,\%\\
R12 & explicit mode, 15 outer iterations, water rising into cells dry at the head of linearization & converged; outflow within 0.6\,\%; MODFLOW error $10^{-7}$\\
R13 & stiff low-storage slope, budget 1, 12 outer iterations, MODFLOW unconverged & neither residual nor shortfall handed to Raven: outflow 1.01 times the iterated run's\\
R15b & a negative-recharge cut on a draining soil zone: 8.67~mm cut while the wet soil drains & soil zone after step 1 at least its field-capacity water (126.08 against 120~mm)\\\bottomrule
\end{longtable}

\section{Regression tests}\label{sec:va1}
Each test of Table~\ref{tab:va1} fails on a mutation of the code or of the check it covers. The full suite
(\cmd{run\_all.py}, full mode) passed 45 of 45 steps; the two USGS example steps were not run (their examples are not on
the test machine) and three permanent skips of \cmd{test\_confined\_wet} (W4, W6, W11) are reported as such.

\begin{longtable}{P{0.19\linewidth}P{0.45\linewidth}P{0.27\linewidth}}
\caption{Regression tests and their results.}\label{tab:va1}\\
\toprule\hdr Test & Case & Result\\\midrule\endfirsthead
\caption[]{Regression tests (continued).}\\
\toprule\hdr Test & Case & Result\\\midrule\endhead
B8 (\cmd{test\_mf6\_budget}) & without \cmd{:SoilZone PHYSICAL}, \cmd{:SolverMaxOuterIterations} 2 (10 unconverged steps) and 1 (every step), and the Liard example (rivers and lakes) at 2: each row's residual over all its columns equals MODFLOW's negative residual on unconverged steps and 0 otherwise; run total ``discharge withheld''; at most ten warnings & step 1: 1\,039.4~m$^3$/d withheld = $-$balance error; run total 1\,062.5~m$^3$ = sum of rows; limit 1: 18\,511~m$^3$; 10 warnings\\
R15b (\cmd{test\_robustness}) & Table~\ref{tab:vszrobust} & 126.08~mm $\ge$ 120~mm\\
I5 (\cmd{test\_vz\_initial}) & steady start for $(1-F)\,q$: hummocks $\pm3$~m, water table 3~m deep ($F=0.16$), 2~mm/d, \cmd{RICHARDS}, \cmd{KINEMATIC}, \cmd{LAG}, $K=1$ and 3 & step-1 outflow = inflow = 100.8~m$^3$, 25.6~m$^3$ rejected (120.0)\\
I6 (\cmd{test\_vz\_initial}) & a column that cannot carry $q$ (\cmd{K\_VERT} 1~mm/d under 2~mm/d) & storage unchanged; outflow = inflow = 40.000~m$^3$, 40.000 rejected; warning\\
C10 (\cmd{test\_vz\_classes}) & hotstart after a step that dropped a class column ($K=3$; after 4, 10 and 13 days) & no mismatch warning; restart within $7.5\times10^{-7}$ of the continuous run\\
I9--I14 (\cmd{test\_inputs}) & Table~\ref{tab:vszinputs} & all as required\\
T3 (\cmd{test\_hillslope}) & drying checks on the HRUs that connect, the toe's recession as reference; \cmd{--selftest} & Table~\ref{tab:vsoilzone1}; four mutations each fail\\
C9 (\cmd{test\_vz\_classes}) & R20 (\cmd{band\_e43}, \cmd{LAG}) with three depth classes through \cmd{test\_rough\_terrain.py} & closure 0.0001 (limit 0.004); 0 failed runs\\
\cmd{run\_all.py} & quick mode reaches every script; skipped checks counted & ``PASS (3 skipped)'' on \cmd{soilzone\_confined\_wet}\\
\cmd{test\_sz\_models} & standard identity limits, no unconverged step; P1 against the manual's numbers; P3 re-forms & worst identity $4\times10^{-10}$; P1 6.06 $\to$ 15.62~d, peak ratio 0.55\\
\cmd{sz\_check.py} & self-test of the column identity: exact books, and a $10^{-4}$~m$^3$ leak per step on a 9\,776~m$^3$ store & 0, and $2.4\times10^{-7}$ (fails)\\
\cmd{cmake\_checks} 15--19 & Table~\ref{tab:vcmake} & 60 parts in 19 checks pass; each of 15--19 fails on its mutation\\\bottomrule
\end{longtable}

\begin{longtable}{P{0.07\linewidth}P{0.60\linewidth}P{0.25\linewidth}}
\caption{CMake checks of the MODFLOW 6 build integration (\cmd{tests/cmake\_checks.py}, offline; checks 1--14: Section~\ref{sec:building}).}\label{tab:vcmake}\\
\toprule\hdr Check & What is checked & Result\\\midrule\endfirsthead
\caption[]{CMake checks of the MODFLOW 6 build integration (continued).}\\
\toprule\hdr Check & What is checked & Result\\\midrule\endhead
15 & offline default: with the source tools but an unreachable git repository the default falls back at configure time to \cmd{PREBUILT} (a local release folder: the library target builds without network) or, with no release either, to \cmd{OFF} with a warning; a reachable repository gives \cmd{SOURCE} & as required\\
16 & default build type: a single-configuration generator without \cmd{CMAKE\_BUILD\_TYPE} configures \cmd{Release} (\cmd{-O3 -DNDEBUG}); a given \cmd{Debug} is kept & as required\\
17 & version floor: 6.4.0 as \cmd{RAVEN\_MF6\_VERSION} (\cmd{SOURCE} and \cmd{PREBUILT}), in a zip's or a source folder's \cmd{meson.build} or an archive name stops the configure naming 6.5.0; 6.7.0 configures with a warning naming 6.8.1 & stops in all five ways; warning\\
18 & runtime bundling, emulated on Linux: macOS with Mach-O stand-ins (\cmd{@rpath} chains resolved through \cmd{LC\_RPATH}, every reference rewritten to \cmd{@loader\_path}, every edited file signed); Windows with stand-in \cmd{objdump} and \cmd{dumpbin} (the DLLs \cmd{libgfortran-5.dll} imports installed) & all libraries installed and rewritten; real macOS and Windows not tested\\
19 & messages: a folder given as \cmd{RAVEN\_MF6\_PATH}; a SHA-256 that is not 64 hex digits & named correctly\\\bottomrule
\end{longtable}
